% !PRINTED 204
\weberpage{204}

Ist $a$ ein $\nu$-zähliger Pol der Gruppe $G$ und

\weq{1}{a = \Theta(a) = \Theta^2(a)\cdots = \Theta^{\nu-1}(a),}

so ist der zweite Pol $a'$ von $\Theta$ nach \S.\,52, 4. gleichfalls $\nu$-zählig,
und es ist

\weq{2}{a' = \Theta(a') = \Theta^2(a')\cdots = \Theta^{\nu-1}(a').}

Giebt es nun in der Gruppe $G$ mehr als ein System conjugirter
$\nu$-zähliger Pole, so können $a, a'$ entweder in demselben
oder auch in verschiedenen dieser Systeme vorkommen.

Giebt es aber nur \gesperrt{ein} System $\nu$-zähliger Pole, so muss $a$
und $a'$ in demselben Systeme vorkommen, und es muss also in $G$
eine nicht unter den Potenzen von $\Theta$ enthaltene Substitution $\psi$
existiren, so dass

\weq{3}{a' = \psi(a)}

ist. Daraus folgt mit Anwendung von (1) und (2):
\[
\psi(a) = \Theta\psi(a), \quad a = \psi^{-1}\Theta\psi(a),
\]
woraus zu schliessen ist, dass $\psi^{-1}\Theta\psi$ unter den Potenzen von
$\Theta$ vorkommt, und dass daher nach (2) auch
\[
a' = \psi^{-1}\Theta\psi(a'), \quad \psi(a') = \Theta\psi(a')
\]
sein muss. Es ist also auch $\psi(a')$ ein Pol von $\Theta$, und weil $\psi(a')$
nicht $= a'$ sein kann, weil sonst $a'$ ein Pol von $\psi$ wäre, was
nicht sein kann, da $\psi$ nicht unter den Potenzen von $\Theta$ vorkommt,
so ist

\weq{4}{\psi(a') = a.}

Wenn man nun durch Transformation der Gruppe die Pole
$a$ und $a'$ nach $0$ und $\infty$ bringt, so erhält $\Theta$ die Form
\[
\Theta(x) = \varepsilon x,
\]
worin $\varepsilon$ eine primitive $\nu^{\text{te}}$ Einheitswurzel ist, und $\psi$ muss wegen
(3) und (4) die Form haben:
\[
\psi(x) = \frac{c}{x},
\]
worin $c$ eine Constante ist. Durch eine abermalige Transformation
der Gruppe kann man der Constanten $c$ jeden beliebigen Werth,
z. B. durch Transformation mit der Substitution $x\sqrt{c}$ für $x$, den
Werth $1$ geben. Wir erhalten also folgenden Satz:

1.~\gesperrt{Wenn in der Gruppe $G$ nur ein System conjugirter
 $\nu$-zähliger Pole vorkommt, so kann man}
%CONT>>
